We compute six pathway metrics from attention maps, each capturing a distinct aspect of routing structure:

\begin{table}[H]
\centering
\small
\begin{tabular}{llp{5cm}}
\toprule
\textbf{Metric} & \textbf{Dimension} & \textbf{Interpretation} \\
\midrule
1. Routing Sparsity & Scalar & Effective number of active attention paths (Rényi entropy) \\
2. Path Competition & Scalar & Dominance of top attention weight: $\max(W) / \text{mean}(W)$ \\
3. Path Efficiency & Scalar & Energy concentration in top-$k$ weights (20\% rule) \\
4. Routing Entropy & Scalar & Shannon entropy of attention distributions \\
5. Inter-head Divergence & Scalar & Mean KL-divergence between attention heads \\
6. Layer Stability & Scalar & Cosine similarity between consecutive layer attentions \\
\bottomrule
\end{tabular}
\caption{Pathway Metrics. All metrics are computed per layer, then averaged.}
\label{tab:metrics}
\end{table}

\subsection*{Motivation}

\begin{enumerate}
    \item \textbf{Routing Sparsity}: Measures effective dimensionality of the path distribution. Sparse attention reduces this.
    \item \textbf{Path Competition}: Quantifies "winner-take-all" dynamics. Higher competition indicates concentrated routing.
    \item \textbf{Path Efficiency}: Captures how much model capacity is utilized; more efficient routing allocates mass to fewer paths.
    \item \textbf{Routing Entropy}: Inverse measure of organization. Lower entropy $\leftrightarrow$ more structured routing.
    \item \textbf{Inter-head Divergence}: Tests whether different attention heads develop specialized roles or are redundant.
    \item \textbf{Layer Stability}: Checks whether routing patterns are consistent across layers (indicating learned structure) or noisy.
\end{enumerate}

Collectively, these metrics form a 6-dimensional pathway signature $\mathcal{P} \in \mathbb{R}^6$.

\subsection*{Normalization}

All metrics are computed on batch-averaged attention maps and then z-scored to unit variance before projection or prediction.
